\section{Related Work}
\label{sec:related}

\paragraph{Attribution graphs and replacement models.}
Attribution graphs operationalize feature$\to$logit pathways in a local replacement model that linearizes the residual stream through CLT or SAE features and freezes attention~\cite{ameisen2025circuittracing,lindsey2025biology}. The Neuronpedia platform exposes these graphs for selected models and displays per-feature cards at scale~\cite{neuronpedia2025,lindsey2025landscape}. \texttt{Circuit-tracer} provides the open-source reference implementation~\cite{circuittracer2025} for computing these graphs. Our work operates entirely \emph{downstream} of attribution: we do not propose a new attribution method, but we do produce supernodes that Neuronpedia tooling can render and pin. Up to our work, the supernode clustering had to be produced by hand.

\paragraph{Causal circuits with transcoders.}
\citet{hannaameisen2026planning} use cross-layer transcoder feature circuits to study whether language models latently plan rhyming words, structuring their analysis around two falsifiability conditions (predict-token-before-emission, causally-affect-token). \citet{marks2024dictionary} apply sparse-feature circuit analysis to subject--verb agreement; we contribute a domain-general harness and a proposal on how to generate groupings of attribution graphs on scale.


\paragraph{Causal interpretability and matched random controls.}
\citet{geiger2025causal} formalize three levels of interpretive claim and emphasize that \emph{distinguishing} levels requires intervention. \citet{heap2025sparse} show that SAE auto-interpretability scores~\cite{bills2023autointerp,paulo2024autointerp} and standard SAE reconstruction metrics can be similar for randomly-initialized and trained transformers, highlighting the need for metrics anchored to causal effect. Our matched-random-control protocol (\S\ref{sec:method:harness}) is designed to address this critique: any structural property a random feature set shares with the labeled one (count, layer histogram, eligibility) is held constant; only the concept alignment changes.
